\documentclass[11pt]{article}
\usepackage[margin=1in]{geometry}
\usepackage{amsmath,amssymb,amsthm,mathrsfs,hyperref}
\newtheorem{lemma}{Lemma}\newtheorem{theorem}{Theorem}
\newtheorem{proposition}{Proposition}
\begin{document}\setlength{\emergencystretch}{1em}
\title{Unit-digit terminal families at ordinary gap six}
\author{Internal proof; no priority or external-review claim}\date{18 November 2026}\maketitle
Let $3\le c\le6$, $r\ge1$, $L=7^r$, $N=L+c$, and let the
raising operator be $\mathcal A=z-zD_z+N+5$.
The six normalized cyclic rows are
\[
 V_a(Q;v)=\sum_{\ell\ge0}
 \frac{7^\ell a!}{(a+6\ell)!}d_{a+6\ell}(Q)v^\ell,\quad
 d_j(Q)=\sum_{\ell=0}^j(-1)^\ell\binom j\ell
 \binom{(N+5-\ell)/6+N-1}{N}Q(N+5-\ell),
\]
for $\deg Q\le4$, on $z^6=7v$.
These rows are obtained from the exact finite identity
\[
 \frac{\Pi_N(\mathcal A)Q(\mathcal A)1}{6^NN!}
       =\sum_{j=0}^{N+4}\frac{d_j(Q)}{j!}z^j,\qquad
 \Pi_N(X)=\prod_{\ell=0}^{N-1}(X+6\ell).
\]
It follows by formal conjugacy and finite backward differences;
no exponential is evaluated at the boundary.
For the moment bridge write
\[
 T_N(t)=T_{N,N,6}(t)=
             \sum_{j=0}^N\binom Nj(N)_{6j}t^j .
\]
They are integral, since the binomial polynomial is integer-valued
on $\mathbb Z_7$ and the cyclic weights have nonnegative valuation.
Write $S=S_r(v)=1+\sum_{i=1}^r(-1)^iv^{(7^i-1)/6}$ and
$\mathscr D=v\partial_v+(S-1)\partial_S$.

\begin{lemma}[Actual uniform seeds]\label{lem:seeds}
The ordinary reduced seeds at $Q=1$ are zero except
\[
 F_c=(-1)^c(1+cS)\quad(c\le5),\qquad
 F_{c+1}=(c+1)(-1)^{c+1}S\quad(c+1\le5).
\]
Every entire row $V_a$, $a<c$, is seven-divisible.
The divided seeds $R_a(1)=\overline{V_a(1)/7}$ are
\[
\begin{array}{c|l}
 c& (R_0,\ldots,R_{c-1})\\\hline
3&(6,\ 1+2S,\ S-2)\\
4&(1,\ 3+2S,\ 2S-3,\ -S-1)\\
5&(6+vS,\ 3-3S,\ 2-2S,\ 2S-2,\ 3S-3)\\
6&(1+v(S-1),\ S(v-1),\ 0,\ 0,\ 0,\ 0).
\end{array}
\]
The complete ordinary rows are generated by
$F_a(XQ)=(c+5-a+\mathscr D)F_a(Q)+aF_{a-1}(Q)$,
with $F_{-1}=0$.
The divided rows satisfy
\[
 R_a(XQ)=(c+5-a+\mathscr D)R_a(Q)+aR_{a-1}(Q),\quad a\ge1,
\]
and
$R_0(XQ)=(c+5+\mathscr D)R_0(Q)+vF_5(Q)$.
\end{lemma}
\begin{proof}
Use the exact diagonal series and integral formal Artin--Hasse
reparametrization in the proofs of \texttt{lem:diagonal} and
\texttt{lem:first-row} of
\href{GAP_SIX_TWO_POWER_PLUS_ONE_SFC.tex}
{the two-power-plus-one note}.
Here $\sum_nV_a^{(n)}(1)t^n=a!\kappa\mathcal C_{6,a}e^{zH}$,
with $R-R^7=t$, $H=1-R$, $\kappa\equiv R-7t\pmod{49}$.
Let $T=\sum t^{7^i}$ and $A_0=\sum a_it^{7^i}$,
$a_0=-1$, $a_i=-1-va_{i-1}^7$.
Then $a_r=-S$ in characteristic seven.

Modulo seven the section is $(1+T)A_0^a$.
Extracting $[x^cy]$ after $T=x+y$, $A_0=-x-Sy$ gives precisely
the two stated ordinary seeds.
All other exponents use a number of selected powers different
from $c+1$ modulo six, or fewer than the digit sum $c+1$.
Raising preserves the initial zero rows $a<c$, proving their
whole-row seven-divisibility.

For $1\le a<c$ and $c\le5$, the divided seed is exactly
\[
 [x^cy]P(x+y)(-x-Sy)^{a-1}
            \{-x-Sy-a(1+x+y)\},\quad
 P(T)=T+3T^2+5T^3+5T^4+3T^5+T^6.
\]
The omitted carry and Artin--Hasse terms have exponents
$a$ or $a+1$ modulo six, not $c+1$.
Only selected-power degree $c+1$ can contribute: the largest
possible degree is $a+6\le c+5<c+7$, and every carry lowers digit
sum by six. Thus there is no intermediate-digit omission.
The divided zero row is $(c+1)[T^{c+1}]P$ for $c\le4$.
For $c=5$ its Artin--Hasse correction is
$6v[x^5y](-x-Sy)^6=vS$.

For $c=6$, the two carry coefficients are
$[t^N]C(T)=1$ and $[t^N]C(A_0)=a_0^6a_r=-S$.
Hence the first divided row acquires $-1+vS$ beyond its
$P$ contribution $1-S$, giving $S(v-1)$.
The zero row is
\[
 1+6v[x^6y](x+y)(-x-Sy)^6=1+v(S-1).
\]
The same polynomial extraction gives zero for $a=2,\ldots,5$.
All statements are formal coefficient identities; no boundary
exponential is evaluated.
The exact row recurrence, before evaluation, is
\[
\begin{aligned}
 V_a(XQ)&=(N+5-a-6vD_v)V_a(Q)+aV_{a-1}(Q),\quad a\ge1,\\
 V_0(XQ)&=(N+5-6vD_v)V_0(Q)+\frac{7v}{5!}V_5(Q).
\end{aligned}
\]
First reduce it to propagate the ordinary rows; their initial zeros
at $a<c$ remain zero for every column $Q=1,X,\ldots,X^4$.
Then divide those entire rows by seven and reduce again.
Use $N+5\equiv c+5$, $-6\equiv1$, $5!\equiv1$, and
$vS_r'=S_r-1$ to get the stated formulas.
\end{proof}

Let $M_c(S,v)$ be the six-by-five matrix formed from $R_a$
for $a<c$, and $F_a$ for $a\ge c$, with columns $1,X,\ldots,X^4$.
It is the actual normalized and divided boundary matrix.
The finite identity and the recurrences prove this exact specialization
at every $r$. Exact controls also compare all thirty entry polynomials
with the integer auxiliary recurrence at $r=1,2$; these are normalization
controls, not an index argument.

\begin{lemma}[Finite necessary constraints]\label{lem:certificate}
For $c=3,4,5$, exact sparse identities in the six maximal minors
give $q_c(S)=0$ and $v=W_c(S)$ at every residual rank drop.
The monic polynomials $q_c$ have degrees $15,20,25$.
In $\mathbb F_7[S]/q_c$, the orbit
$x_0=1$, $x_{j+1}=1-W_cx_j^7$ has periods dividing $14,28,42$,
respectively.
For $c=3,4$, every target $x_j-S$ is a unit in the full period.
For $c=5$ its gcd with $q_5$ is one except when $j=1$ modulo6,
when it is $S-4$; also $W_5(4)=4$ and $q_5'(4)=5$.

For $c=6$, exact maximal-minor identities give $S^6=0$ and
$f(S,v)=0$, with $f(0,v)=(v-1)^5$.
\end{lemma}
\begin{proof}
The complete finite data are retained in
\href{gap_six_small_offset_families_certificate.json}
{the coefficient certificate}.
It contains all four matrices, their six maximal minors, both
necessary-relation polynomials, and explicit multiplier polynomials.
For $c=3,4,5$ the identities are
\[
 q_c(S)=\sum_iH_{c,i}(S,v)\Delta_{c,i}(S,v),\qquad
 v-W_c(S)=\sum_iK_{c,i}(S,v)\Delta_{c,i}(S,v).
\]
For $c=6$ the two targets are $S^6$ and the displayed $f$.
The multiplier counts are $49,103,153,17$.
Direct sparse expansion checks all eight identities.
The verifier rebuilds each matrix from the seeds and raising formulas,
and recomputes every maximal minor, so no elimination algorithm is
trusted as an input.

The three orbit records contain every state, target gcd and return
to1. Exact polynomial arithmetic verifies every entry.
The return to the initial state proves periodicity for arbitrary$r$.
In the $c=5$ record the only nonunit target gcd is $S-4$, exactly
at depths1 modulo6. Evaluation gives $q_5'(4)=5$, $W_5(4)=4$.
The last necessary relation gives $f(0,v)=(v-1)^5$ by exact
coefficient comparison. All assertions hold over the full residue
field, not only its seven prime-field elements.
\end{proof}

\begin{theorem}\label{thm:families}
Ordinary gap-six Strong Factorial holds at every terminal
$N=7^r+c$, $c=3,4,5$, $r\ge0$.
\end{theorem}
\begin{proof}
At $r=0$, $N=4,5,6$ are in the native Eisenstein family,
\texttt{cor:pair} of the
\href{PHYSICAL_PRIME_TOTAL_DEGREE_REDUCTION.tex}{physical-prime note}.
Let $r\ge1$.
That note's \texttt{thm:radius} confines any multiple root of $T_N$
to $u=-1/t=7v$, with $v$ a unit.
The \href{gap_three_hurwitz_structure.tex}{critical-ideal lemma},
\texttt{lem:critical-ideal}, makes the corresponding cyclic matrix
of $B_N,\ldots,B_{N+4}$ lose rank.
This lemma is applied over $\mathbb C$ at an algebraic multiple root;
its rational polynomial minor vanishing then transports to a7-adic
embedding. No7-adic exponential is evaluated.

For the fixed shape $N$,
\[
 B_{N+j}=\Pi_N(\mathcal A)
       \prod_{\ell=N}^{N+j-1}(\mathcal A+6\ell)1.
\]
The five multiplier polynomials are monic of consecutive degrees,
so their columns are a unit triangular change from $Q=1,X,\ldots,X^4$.
The exact finite identity, nonzero constant normalizations, and global
row divisions preserve characteristic-zero rank.
All resulting matrix entries are integral.
Normalize a nonzero kernel vector by one coefficient of minimum
valuation. Its reduction is nonzero and would give a kernel of
$M_c(S_r(\bar v),\bar v)$ over$k$.

For $c=3,4$, the unit orbit targets in Lemma~\ref{lem:certificate}
exclude this: evaluation sends $x_r$ to$S_r(\bar v)$ and cannot
send a unit to zero.
For $c=5$, the only possible residual point is $S=4,v=4$, and
only the class $r=1$ modulo6 can contain it.
In that class, lift the $q_5$ multiplier identity integrally and
substitute $S=S_r(v)$ in the lifted multipliers.
Applying them to the actual saturated minors gives
$g_r(v)\in\mathbb Z_7[v]$, vanishing at every actual rank drop,
with $\bar g_r=q_5(S_r(v))$.
At $v=4$ its derivative is the unit
\[
 q_5'(4)S_r'(4)=5(4-1)/4=5\cdot6=2\quad\text{in }\mathbb F_7.
\]
For any distinct $x,y$ reducing to4, the polynomial divided
difference $(g_r(x)-g_r(y))/(x-y)$ reduces to this unit.
Thus $g_r$ is injective on the entire residue disk
$4+\mathfrak m$, including arbitrarily ramified extensions.
There is at most one actual rank-drop point.

Every algebraic conjugate of a multiple root of the rational
polynomial $T_N$ enters that same disk and the same fixed matrix.
Separability would force its irreducible factor over $\mathbb Q$
to have degree one, hence its root real.
But $N=7^r+5$ is even and, for real$t$,
\[
 T_N(t)=\frac1{\Gamma(N)}
       \int_0^\infty x^{N-1}e^{-x}(1+tx^6)^N\,dx>0.
\]
This contradiction excludes it; the integrand is nonnegative and
positive on an interval.
Therefore all three cases give squarefreeness of$T_N$.

Finally the normalized consecutive moments are
$T_N+(6/N)tT_N'$ and $T_N-(1/N)tT_N'$.
Their determinant is $-7/N\ne0$ and $T_N(0)=1$.
A common zero would be a multiple root of$T_N$.
Degenerate nonzero monomials have no zero positive moment.
\end{proof}

\begin{proposition}[An even-depth subfamily]\label{prop:even-six}
The same SFC assertion holds at $N=7^{2j}+6$, $j\ge0$.
For odd$r$, the first divided reduction leaves only the residue
point $S=0,v=1$; it does not exclude an actual multiple root.
\end{proposition}
\begin{proof}
At a residual rank drop the last part of
Lemma~\ref{lem:certificate} forces $S=0,v=1$ over every residue
field. The actual recursion is $S_0(1)=1$,
$S_{r+1}(1)=1-S_r(1)$, so $S_r(1)=1$ for even$r$.
Thus even$r$ gives full residual rank and the same bridge applies.
The case $r=0$, $N=7$, is native Eisenstein.
For odd$r$, $S_r(1)=0$ is compatible; no further exclusion is
claimed.
\end{proof}

\begin{corollary}[A combined small-offset family]\label{cor:small-offsets}
Ordinary gap-six SFC holds at every terminal $N=7^r+c$,
$2\le c\le5$, $r\ge0$, and at $N=7^{2j}+6$, $j\ge0$.
\end{corollary}
\begin{proof}
Theorem~\ref{thm:families} and Proposition~\ref{prop:even-six}
supply $c=3,4,5$ and the even-depth $c=6$ case.
For $c=2,r\ge1$ use \texttt{thm:family} of the
\href{GAP_SIX_POWER_PLUS_TWO_SFC.tex}{power-plus-two note}.
Its endpoint $r=0,N=3$ is native Eisenstein, \texttt{cor:pair}
of the physical-prime note.
\end{proof}

\begin{corollary}[The refined degree-seven frontier]
\label{cor:degree-seven-frontier}
Put
\[
\begin{split}
 \mathcal U={}&\{b7^s:1\le b\le6,\ s\ge0\}
 \cup\{b7^r+1:1\le b\le6,\ r\ge0\}\\
 &{}\cup\{7^r+c:2\le c\le5,\ r\ge0\}
 \cup\{7^{2j}+6:j\ge0\}.
\end{split}
\]
Ordinary binomial SFC through degree seven is equivalent to
$\gcd_{\mathbb Q[t]}(P_{N-1},P_N)=1$ at every terminal $N\ge6$
outside $\mathcal U$, where
$P_j(t)=\sum_k\binom jk(j+1)_{6k}t^k$.
Any hypothetical common root on those remaining tests has
$v_7(-1/t)=1$ at every algebraic embedding.
\end{corollary}
\begin{proof}
Use \texttt{cor:degree-seven-frontier} of the
\href{GAP_SIX_SCALED_POWER_PLUS_ONE_SFC.tex}{scaled power-plus-one note}.
Corollary~\ref{cor:small-offsets} supplies every additional omitted
terminal. Removing already proved tests preserves sufficiency;
necessity is immediate from inclusion among the degree-seven tests.
\end{proof}
The still-compatible odd-depth $N=7^r+6$ residue point is a
limitation of this first divided reduction, not an SFC counterexample.
The remaining terminals, all-index gap-six and degree-seven SFC,
and unrestricted FC/SFC remain open.
No all-positive-Gamma-shape extension is asserted here.
Written formal-series and rank arguments supply the unbounded depth;
exact finite orbits and literal profiles have their distinct stated roles.
Internal audit and replay do not imply independent whole-proof replay,
formal verification, external review or priority.
\end{document}
